% META: {"id": "1SPE-SECDEG-CO-029", "chapitre": "1SPE-SECOND-DEGRE", "type_objet": "corrige", "exercice_ref": "1SPE-SECDEG-EX-029", "capacites_codes": ["C3", "C5"], "sources_inspiration": [], "status": "generated"}
% BEGIN-VERIFY
% from sympy import symbols, solve, Rational, sqrt
% x = symbols('x')
% f = x**2 - 5*x + 3
% delta_f = (-5)**2 - 4*1*3
% assert delta_f == 13
% sols_f = sorted(solve(f, x))
% assert sols_f == [(5 - sqrt(13))/2, (5 + sqrt(13))/2]
% # f <= 0 sur [x1, x2]
% assert f.subs(x, 0) > 0   # hors
% # g(x) = x^2 + 4x + 5
% g = x**2 + 4*x + 5
% delta_g = 4**2 - 4*1*5
% assert delta_g == -4   # delta < 0 et a > 0 => toujours positif
% assert g.subs(x, 0) > 0
% assert g.subs(x, -2) > 0
% assert g.subs(x, -100) > 0
% END-VERIFY
\begin{corrige}{1SPE-SECDEG-EX-029}

\textbf{1. Résolution de $x^2 - 5x + 3 = 0$.}

\commentaireMarge{$a=1$, $b=-5$, $c=3$ ; solutions exactes avec $\sqrt{\Delta}$.}

$a = 1$, $b = -5$, $c = 3$.
\[\Delta = (-5)^2 - 4 \times 1 \times 3 = 25 - 12 = 13.\]

$\Delta = 13 > 0$ : l'équation admet deux solutions réelles distinctes.

\[x_1 = \frac{5 - \sqrt{13}}{2} \qquad \text{et} \qquad x_2 = \frac{5 + \sqrt{13}}{2}.\]

\medskip
\textbf{2. Résolution de $x^2 - 5x + 3 \leqslant 0$.}

\commentaireMarge{$a = 1 > 0$ : $f \leqslant 0$ entre les racines.}

Comme $a = 1 > 0$, le trinôme est négatif ou nul entre ses racines.

L'ensemble des solutions est :
\[\mathcal{S} = \left[\frac{5 - \sqrt{13}}{2}\,;\,\frac{5 + \sqrt{13}}{2}\right].\]

\medskip
\textbf{3. Discriminant de $g(x) = x^2 + 4x + 5$.}

\commentaireMarge{Calculer $\Delta_g = b^2 - 4ac$.}

$a = 1$, $b = 4$, $c = 5$.
\[\Delta = 4^2 - 4 \times 1 \times 5 = 16 - 20 = -4.\]

$\Delta = -4 < 0$.

\medskip
\textbf{4. Signe de $g(x)$ et résolution de $g(x) > 0$.}

\commentaireMarge{$\Delta < 0$ et $a > 0$ : $g$ est strictement positive sur $\mathbb{R}$.}

Comme $\Delta < 0$ et $a = 1 > 0$, le trinôme $g$ n'a pas de racine réelle et son coefficient dominant est positif.

Donc $g(x) > 0$ pour tout $x \in \mathbb{R}$.

L'ensemble des solutions de $g(x) > 0$ est $\mathcal{S} = \mathbb{R}$.

\baremeIndicatif{Q1 : 2 pts (calculer) ; Q2 : 2 pts (raisonner) ; Q3 : 1 pt (calculer) ; Q4 : 2 pts (raisonner, communiquer)}
\end{corrige}
